\chapter{Програмна реалізація}

У цьому розділі представлено програмну реалізацію розглянутих у роботі
стеганографічних методів та їхній порівняльний аналіз за критеріями якості
контейнера й стійкості до стегоаналізу. Реалізація виконана мовою
програмування Python з використанням бібліотек \texttt{NumPy} (масивні
обчислення), \texttt{Pillow} (робота із зображеннями), \texttt{SciPy}
(дискретне косинусне перетворення та статистичні розподіли). У першому
підрозділі досліджуються методи для растрових контейнерів формату
\texttt{.bmp} (стандартний LSB та удосконалений LSB-зіставлення), у
другому~--- для контейнерів формату \texttt{.jpeg} (стандартний DCT та
алгоритм Shield).

Для кожної пари «стандартний метод~--- удосконалений метод» проведено
експерименти з різним обсягом приховуваного повідомлення, обчислено пікове
співвідношення «сигнал/шум» (PSNR) та середньоквадратичну похибку (MSE), а
також застосовано методи стегоаналізу (хі-квадрат тест та RS-аналіз) для
оцінки стійкості до виявлення.

\section{Приховування інформації у файлах формату .bmp}

Формат \texttt{.bmp} без стиснення є оптимальним вибором для методів
просторової області, оскільки гарантує точне збереження кожного пікселя у
файлі без втрат, що унеможливлює пошкодження прихованого повідомлення.
В якості контейнера використано зображення розміром $1280 \times 853$
пікселів у режимі RGB, що забезпечує максимальну ємність вбудовування
$1280 \cdot 853 \cdot 3 / 8 = 409\,440$ байтів за умови використання трьох
колірних каналів.

\subsubsection*{Стандартний метод LSB}

Метод заміни найменш значущого біта реалізує послідовну модифікацію
останнього біта кольорових компонент пікселів зображення-контейнера. Для
підвищення безпеки позиції вбудовування обираються псевдовипадково на
основі секретного ключа $K$. Алгоритм вбудовування реалізовано наступним
чином:
\begin{itemize}
    \item Завантажуємо контейнерне зображення формату \texttt{.bmp} та
    перетворюємо його у масив пікселів розмірності $H \times W \times 3$.
    \item Секретне повідомлення кодуємо у двійкову послідовність
    $b_1 b_2 \dots b_L$, де кожен символ подається у вигляді 8-бітного коду.
    \item Ініціалізуємо псевдовипадковий генератор ключем $K$ та формуємо
    послідовність позицій $(y_i, x_i, c_i)$, $i = 1, \dots, L$, де $y_i$,
    $x_i$~--- координати пікселя, $c_i \in \{R, G, B\}$~--- індекс
    колірного каналу.
    \item Для кожної позиції замінюємо найменш значущий біт відповідної
    компоненти: $S_{y_i, x_i}^{(c_i)} := (C_{y_i, x_i}^{(c_i)} \,\&\, \texttt{0xFE}) \lor b_i$.
    \item Зберігаємо модифіковане зображення у форматі \texttt{.bmp}.
\end{itemize}
Видобування повідомлення виконується у зворотному порядку: за тим самим
ключем $K$ відновлюємо послідовність позицій, зчитуємо молодший біт кожної
зазначеної компоненти та групуємо біти по 8 для відновлення вихідного
тексту.

% Рис. 3.1 — місце для зображення (порожній + заповнений LSB)
\begin{figure}[h!]
    \centering
    % \includegraphics[width=0.8\textwidth]{lsb_comparison.png}
    \rule{0.8\textwidth}{4cm}
    \caption{Різниця між порожнім і заповненим контейнером (LSB)}
    \label{fig:lsb-diff}
\end{figure}

\subsubsection*{Метод LSB-зіставлення (LSB Matching)}

На відміну від стандартного LSB, метод LSB-зіставлення при невідповідності
біта повідомлення молодшому біту пікселя не замінює його, а випадково
змінює значення пікселя на одиницю в більший або менший бік. Завдяки цьому
не виникає характерного для класичного LSB ефекту спарювання частот
$n_{2k} \approx n_{2k+1}$, який детектує хі-квадрат стегоаналіз. Алгоритм
вбудовування виконується таким чином:
\begin{itemize}
    \item Завантажуємо контейнер та конвертуємо повідомлення у двійкову
    послідовність аналогічно стандартному методу LSB.
    \item За ключем $K$ генеруємо послідовність позицій
    $(y_i, x_i, c_i)$ для вбудовування.
    \item Для кожної позиції перевіряємо умову збігу:
    \begin{itemize}
        \item якщо $\mathrm{LSB}(C_{y_i, x_i}^{(c_i)}) = b_i$, піксель
        залишається без змін;
        \item інакше випадково обираємо $r \in \{-1, +1\}$ з рівною
        ймовірністю $1/2$ та обчислюємо нове значення
        $S_{y_i, x_i}^{(c_i)} = C_{y_i, x_i}^{(c_i)} + r$.
    \end{itemize}
    \item На краях діапазону робимо корекцію: якщо $C = 0$, то $r := +1$;
    якщо $C = 255$, то $r := -1$, щоб уникнути виходу за межі $[0, 255]$.
    \item Зберігаємо модифіковане зображення.
\end{itemize}
Видобування повідомлення повністю ідентичне стандартному методу LSB,
оскільки після перетворення $C \pm 1$ молодший біт результату завжди
дорівнює $b_i$.

\begin{figure}[h!]
    \centering
    % \includegraphics[width=0.8\textwidth]{lsbm_comparison.png}
    \rule{0.8\textwidth}{4cm}
    \caption{Різниця між порожнім і заповненим контейнером (LSB-зіставлення)}
    \label{fig:lsbm-diff}
\end{figure}

\subsubsection*{Стегоаналіз та порівняння методів}

Для оцінки ефективності та стійкості до стегоаналізу обчислено PSNR і MSE
між контейнером та стегозображенням, а також застосовано два методи
стегоаналізу: хі-квадрат тест Вестфельда--Пфіцмана (оцінка вирівнювання
частот пар значень) та RS-аналіз Фрідріх (оцінка відсотка модифікованих
пікселів). Результати наведено в таблиці~\ref{tab:bmp-results}.

\begin{table}[h!]
\centering
\caption{Результати експериментів для методів LSB та LSB-зіставлення
($1280 \times 853$, ємність $409{,}4$\,Кбайт)}
\label{tab:bmp-results}
\small
\begin{tabular}{|c|c|c|c|c|c|c|c|c|}
\hline
\multirow{2}{*}{\shortstack{Розмір\\(байт)}} &
\multicolumn{4}{c|}{LSB (стандартний)} &
\multicolumn{4}{c|}{LSB-зіставлення} \\
\cline{2-9}
& PSNR & MSE & $p_{\chi^2}$ & RS, \% & PSNR & MSE & $p_{\chi^2}$ & RS, \% \\
\hline
$0$ (контейнер)   & ---   & ---    & 0.000 & 0.27 & ---   & ---    & 0.000 & 0.27 \\
$1\,000$          & 77.25 & 0.0012 & 0.000 & 0.20 & 77.25 & 0.0012 & 0.000 & 0.28 \\
$50\,000$         & 60.29 & 0.0609 & 0.000 & 0.83 & 60.29 & 0.0609 & 0.000 & 0.17 \\
$100\,000$        & 57.27 & 0.1219 & 0.000 & 1.84 & 57.27 & 0.1219 & 0.000 & 0.16 \\
$200\,000$        & 54.25 & 0.2442 & 0.000 & 2.98 & 54.25 & 0.2442 & 0.000 & 0.12 \\
$400\,000$        & 51.24 & 0.4886 & 0.100 & 3.45 & 51.24 & 0.4886 & 0.000 & 0.19 \\
\hline
\end{tabular}
\end{table}

Аналіз отриманих даних показує таке. По-перше, обидва методи демонструють
ідентичні значення PSNR та MSE, що очікувано: в обох випадках при
неспівпаді біта значення пікселя змінюється на одиницю (у LSB~--- через
заміну біта 0$\leftrightarrow$1, у LSB-зіставленні~--- через $\pm 1$), а
середня кількість модифікацій дорівнює $L/2$ на повідомлення довжиною $L$
бітів. По-друге, RS-аналіз чітко демонструє різну поведінку методів:
оцінка частки модифікованих пікселів для стандартного LSB монотонно зростає
з~$0.27\%$ (порожній контейнер) до $3.45\%$ (повідомлення розміром
$400$~Кбайт), тоді як для методу LSB-зіставлення вона залишається
статистично нерозрізнюваною від показника порожнього контейнера ($0.12 \dots
0.28\%$). По-третє, хі-квадрат тест починає реагувати на стандартний LSB
лише при максимальному заповненні ($p_{\chi^2} = 0.1$ при 98\% ємності),
тоді як для LSB-зіставлення $p$-значення залишається нульовим для всіх
обсягів повідомлення. Це підтверджує теоретичне передбачення: модифікація
$\pm 1$ не створює ефекту спарювання частот пар PoV, який лежить в основі
обох тестів.

\section{Приховування інформації у файлах формату .jpeg}

Формат \texttt{.jpeg} є природним носієм для методів частотної області,
оскільки сам алгоритм JPEG-стиснення передбачає блочне дискретне косинусне
перетворення з подальшим квантуванням за стандартною таблицею. Це дозволяє
вбудовувати інформацію безпосередньо у квантовані DCT-коефіцієнти без
додаткових перетворень. У якості контейнера використано зображення розміром
$1280 \times 853$ пікселів, з якого виділено канал яскравості $Y$ кольорового
простору $\mathrm{YCbCr}$. Реалізація використовує стандартну матрицю
квантування JPEG з коефіцієнтом якості $\alpha = 50$ (формула (2.7)).

\subsubsection*{Стандартний метод DCT}

Метод реалізує вбудовування шляхом заміни молодшого біта квантованих
середньочастотних коефіцієнтів DCT, обраних псевдовипадково за ключем.
Алгоритм виконується таким чином:
\begin{itemize}
    \item Завантажуємо зображення та виділяємо канал яскравості $Y$;
    центруємо значення відніманням $128$.
    \item Розбиваємо канал на неперекривні блоки розміром $8 \times 8$
    пікселів.
    \item До кожного блоку застосовуємо дискретне косинусне перетворення
    (формула (2.1)) та квантуємо отримані коефіцієнти діленням на матрицю
    $Q$ з округленням до найближчого цілого (формула (2.6)).
    \item За ключем $K$ генеруємо псевдовипадковий порядок обходу блоків та
    порядок позицій усередині блоку, обмежуючись середньочастотною областю
    (індекси 6--27 у зигзаг-скануванні).
    \item Для кожної обраної позиції замінюємо молодший біт абсолютного
    значення квантованого коефіцієнта на наступний біт повідомлення.
    \item Виконуємо обернене квантування та обернене DCT для кожного блоку;
    реконструйовані блоки об'єднуємо у вихідне зображення.
\end{itemize}
Видобування повідомлення виконується аналогічним повторенням етапів
розбиття, DCT та квантування з подальшим читанням молодших бітів коефіцієнтів
у тому самому порядку.

\begin{figure}[h!]
    \centering
    % \includegraphics[width=0.8\textwidth]{dct_comparison.png}
    \rule{0.8\textwidth}{4cm}
    \caption{Різниця між порожнім і заповненим контейнером (DCT)}
    \label{fig:dct-diff}
\end{figure}

\subsubsection*{Алгоритм Shield}

Алгоритм Shield усуває основну вразливість стандартного DCT~--- модифікацію
коефіцієнтів зі значеннями $0$ та $\pm 1$, що деформує характерну
лапласову форму гістограми DCT-коефіцієнтів. У Shield такі коефіцієнти
повністю виключаються з простору вбудовування, а кількість вбудовуваних
бітів адаптивно залежить від абсолютної величини коефіцієнта згідно з
таблицею~2.1. Алгоритм виконується таким чином:
\begin{itemize}
    \item Виконуємо ті самі підготовчі етапи, що й для стандартного DCT:
    виділення $Y$-каналу, розбиття на блоки $8 \times 8$, DCT та
    квантування.
    \item За ключем $K$ генеруємо псевдовипадковий порядок обходу блоків.
    \item У кожному блоці послідовно проходимо AC-коефіцієнти у зигзаг-порядку;
    для кожного коефіцієнта $F_Q(u, v)$ обчислюємо ємність $n$ за таблицею
    2.1: $n = 0$ при $|F_Q| \le 1$ (коефіцієнт пропускається), $n = 1$ при
    $|F_Q| \in [2, 8]$, $n = 2$ при $|F_Q| \in [9, 16]$, і так далі.
    \item Якщо $n > 0$, зчитуємо наступні $n$ бітів повідомлення та
    замінюємо ними $n$ молодших бітів абсолютного значення $|F_Q(u, v)|$;
    знак коефіцієнта зберігається.
    \item Виконуємо обернене квантування, обернене DCT та об'єднуємо блоки
    у стегозображення.
\end{itemize}
Видобування здійснюється у дзеркальному порядку: для кожного коефіцієнта,
що задовольняє умову $|F_Q(u, v)| \ge 2$, визначається $n$ за тією ж
таблицею і зчитуються $n$ молодших бітів його абсолютного значення.

\begin{figure}[h!]
    \centering
    % \includegraphics[width=0.8\textwidth]{shield_comparison.png}
    \rule{0.8\textwidth}{4cm}
    \caption{Різниця між порожнім і заповненим контейнером (Shield)}
    \label{fig:shield-diff}
\end{figure}

\subsubsection*{Стегоаналіз та порівняння методів}

Для методів частотної області RS-аналіз непридатний, оскільки його
дискримінаційна функція оперує яскравостями сусідніх пікселів у просторовій
області. Тому як основний інструмент стегоаналізу використовується
хі-квадрат тест, адаптований під DCT-коефіцієнти. Тестування виконується на
парі значень $(0, 1)$ серед середньочастотних квантованих коефіцієнтів, де
вбудовування DCT найсильніше викривлює природне співвідношення (нулі
перетворюються на одиниці). Як показник дисбалансу пари використано
нормовану різницю
$$
\delta_{01} = \frac{|n_0 - n_1|}{n_0 + n_1},
$$
де $n_0$, $n_1$~--- кількості коефіцієнтів зі значеннями $0$ і $\pm 1$
відповідно. Натуральне зображення характеризується високим значенням
$\delta_{01}$ (близьким до $0.76$ для досліджуваного контейнера); вбудовування,
що модифікує пару $(0, 1)$, знижує цей показник. Результати наведено в
таблиці~\ref{tab:jpeg-results}.

\begin{table}[h!]
\centering
\caption{Результати експериментів для методів DCT та Shield
($1280 \times 853$, $\alpha = 50$)}
\label{tab:jpeg-results}
\small
\begin{tabular}{|c|c|c|c|c|c|c|c|c|}
\hline
\multirow{2}{*}{\shortstack{Розмір\\(байт)}} &
\multicolumn{4}{c|}{DCT (стандартний)} &
\multicolumn{4}{c|}{Shield} \\
\cline{2-9}
& PSNR & MSE & $\chi^2$ & $\delta_{01}$ & PSNR & MSE & $\chi^2$ & $\delta_{01}$ \\
\hline
$0$ (контейнер) & ---   & ---     & 205\,929 & 0.7628 & ---   & ---     & 205\,929 & 0.7628 \\
$100$           & 43.72 & 2.758   & 205\,014 & 0.7611 & 43.55 & 2.870   & 205\,978 & 0.7628 \\
$500$           & 41.87 & 4.225   & 201\,655 & 0.7548 & 41.46 & 4.648   & 205\,973 & 0.7628 \\
$2\,000$        & 38.25 & 9.737   & 189\,004 & 0.7307 & 37.56 & 11.401  & 205\,956 & 0.7628 \\
$5\,000$        & 35.03 & 20.436  & 164\,551 & 0.6818 & 34.18 & 24.834  & 205\,928 & 0.7627 \\
$10\,000$       & 32.29 & 38.393  & 126\,943 & 0.5989 & 33.74 & 27.497  & 205\,923 & 0.7627 \\
\hline
\end{tabular}
\end{table}

Аналіз результатів виявляє разючу різницю між методами щодо стійкості до
хі-квадрат стегоаналізу. Для стандартного методу DCT значення $\chi^2$
монотонно спадає з $205\,929$ (порожній контейнер) до $126\,943$ при
вбудовуванні $10\,000$ байт, тобто на $38\%$; коефіцієнт дисбалансу
$\delta_{01}$ зменшується з $0.76$ до $0.60$. Це означає, що пара значень
$(0, 1)$ поступово вирівнюється внаслідок того, що ~50\% коефіцієнтів зі
значенням $0$ змінюють своє значення на $\pm 1$ при заміні їх молодшого
біта. Для алгоритму Shield ситуація принципово інша: значення $\chi^2$ та
$\delta_{01}$ залишаються практично незмінними ($\delta_{01} = 0.7627
\dots 0.7628$) для будь-якого обсягу вбудовуваного повідомлення, оскільки
коефіцієнти зі значеннями $0$ та $\pm 1$ алгоритм пропускає за конструкцією.
Таким чином, хі-квадрат тест на парі $(0, 1)$ не виявляє жодних слідів
вбудовування у стегоконтейнерах, створених алгоритмом Shield.

Окремо варто відзначити перевагу Shield за показником PSNR: при однаковому
обсязі повідомлення $10\,000$ байт алгоритм забезпечує $33.74$\,дБ проти
$32.29$\,дБ для стандартного DCT, що відповідає меншому викривленню
візуальної якості. Це досягається завдяки адаптивній ємності: великі за
модулем коефіцієнти приймають більше бітів інформації, але відносна зміна
коефіцієнта залишається невеликою.

\subsubsection*{Узагальнення результатів}

Проведені експерименти підтверджують теоретичні передбачення розділу 2: як
у просторовій, так і в частотній області удосконалені методи (LSB-зіставлення
та Shield) забезпечують значно вищу стійкість до стандартних методів
стегоаналізу при збереженні (а у випадку Shield~--- покращенні) якості
стегозображення. Стандартні методи LSB та DCT детектуються відповідно
RS-аналізом і хі-квадрат тестом вже при відносно невеликих обсягах
вбудовуваного повідомлення, тоді як їхні удосконалення залишаються
статистично нерозрізнюваними від порожнього контейнера в усьому
досліджуваному діапазоні.
